\begin{frame}{전체 그림: 4주 파이프라인}
  \begin{columns}[T]
    \begin{column}{0.52\textwidth}
      \begin{itemize}
        \item 한 장으로: \textbf{데이터 $\to$ 파이프라인 $\to$ 평가
          $\to$ 발표} (8.1과 같은 4박스, 박스에 쓴 장이 Ch09~15로 변경)
        \item \textbf{붉은 점선}(평가$\to$파이프라인)이 핵심:
          test에서 \kb{누수}(Ch06.3)를 발견하거나, train/val 곡선이
          \kb{과적합의 모양}이면(Ch09.3, 10.2) \textbf{돌아가서 다시}
        \item 구조적 차이: ``정형(Ch07 GBDT)으로 풀 것인가, 비정형
          (Ch09~15 신경망)으로 풀 것인가'' \textbf{분기}가 Box①과 Box②
          사이에 들어옴
        \item 분기의 결과는 뒤의 모든 선택을 결정 $\Rightarrow$
          \textbf{M1(문제 정의)에서 이미 정해야 한다}
      \end{itemize}
    \end{column}
    \begin{column}{0.46\textwidth}
      \centering
      \includegraphics[width=\linewidth]{ch16_1_project_pipeline.png}
      \vskip0.4em
      {\footnotesize 파이프라인 한 장 — M1(정형이면 GBDT? 분기 포함)
      $\to$ 3분할·전처리·신경망 $\to$ test 한 번·수식적 정당화
      $\to$ 발표·리뷰}
    \end{column}
  \end{columns}
  \vskip0.4em
  {\footnotesize ``한 번 지나가면 돌아올 수 없는 것''은
  \textbf{test 한 번의 측정}뿐 — 그 외 모든 단계는 되돌아갈 수 있다.}
\end{frame}
